% Generated draft: requires mathematical and source review.
\documentclass[11pt]{article}
\usepackage[margin=25mm]{geometry}
\usepackage{fontspec}
\setmainfont{Latin Modern Roman}
\usepackage{amsmath,amssymb,mathtools,unicode-math}
\setmathfont{Latin Modern Math}
\usepackage{enumitem,needspace,xurl,hyperref}
\setlength{\parindent}{0pt}
\setlength{\parskip}{5pt}
\setlength{\emergencystretch}{3em}
\newcommand{\N}{\mathbb{N}}
\newcommand{\Z}{\mathbb{Z}}
\newcommand{\Q}{\mathbb{Q}}
\newcommand{\R}{\mathbb{R}}
\newcommand{\C}{\mathbb{C}}
\newcommand{\F}{\mathbb{F}}
\newcommand{\A}{\mathbb{A}}
\newcommand{\PP}{\mathbb{P}}
\newcommand{\E}{\mathbb{E}}
\newcommand{\eps}{\varepsilon}
\newcommand{\dd}{\,\mathrm{d}}
\begin{document}
\section*{UnsolvedMath ID: 30004345}
\textbf{Source problem code:} catalogue entry
\subsection{Definitions and mathematical statement}
The supplied record concerns three modular elliptic curves and a conjectural \(p\)-adic Gross--Kudla--Zagier formula for modified diagonal cycles. Here \(p\) denotes a prime, but the record supplies neither its restrictions nor the fields of definition and reduction hypotheses for the curves. It denotes the modified diagonal cycle by \(\Delta^\dagger_\xi\) and the purported canonical \(p\)-adic height pairing by \(\langle\ ,\ \rangle_{\mathrm{Nek}}\); the cycle, the datum \(\xi\), the pairing's domain, and the choices needed to define it are not specified.

The source asserts the following schematic identity for every admissible pair \(h,h'\):
\[
\begin{aligned}
&\bigl\langle h^*\Delta^\dagger_\xi,
                 h'^*\Delta^\dagger_\xi\bigr\rangle_{\mathrm{Nek}}
 (\sqrt{3}-1)\,
 E_p\!\left(\operatorname{Fil}^+V_pE\right)
\\
&\qquad = L'_p(E,1)\,
 I^\natural(h\otimes h')\,
 \bigl(\text{a power of }2\bigr)
 \left(\frac{\zeta(2)}{\pi^2}\right)^2.
\end{aligned}
\]
Here \(\zeta\) is the Riemann zeta function. The notation \(L'_p(E,1)\) indicates a derivative of a \(p\)-adic \(L\)-function at \(1\), but its construction, differentiation parameter, and normalization are absent. The record does not define \(E\), \(V_pE\), \(\operatorname{Fil}^+\), \(E_p\), or \(I^\natural\), or explain how \(E\) encodes the three curves. It also does not define admissibility, the operations \(h^*,h'^*\), or the tensor \(h\otimes h'\).

The factor \(\sqrt{3}-1\) is retained exactly as supplied; its interpretation in the coefficient field of this \(p\)-adic identity is unspecified. Likewise, the exponent in ``a power of \(2\)'' and its possible dependence on the data are unspecified. Thus the display preserves the supplied conjectural assertion and its universal quantifier, but is not a fully determined mathematical conjecture without additional hypotheses, definitions, and normalizations.
\subsection{Short English statement}
The supplied conjecture relates \(p\)-adic heights of modified diagonal cycles associated with three modular elliptic curves to a derivative of a \(p\)-adic \(L\)-function and additional factors, for all admissible \(h,h'\). The record leaves the admissible data, several objects, and numerical normalizations undefined, so the displayed formula remains schematic. Its supplied open status and literature assessment require independent review.
\subsection{Sources}
\begin{itemize}
\item Local UnsolvedMath registry, permanent ID \(30004345\), source problem code \texttt{catalogue entry}, title ``p-Adic Gross--Kudla--Zagier Formula for Diagonal Cycles,'' in the Oberwolfach Reports Open Problems collection. The supplied record gives proposed year 2019 and no named proposer.
\item Supplied primary-source citation: \emph{Modular Forms (2020)}, Oberwolfach Reports, pp.~3529--3606, DOI \texttt{10.4171/owr/2019/57}. \url{https://doi.org/10.4171/owr/2019/57}. The background describes Oberwolfach Reports as documenting workshops and their open-problem discussions; no exact conjecture number or page for this formula is supplied.
\item Supplied literature reference: Bertolini--Seveso--Venerucci, ``diagonal cycles and Euler systems I.'' \url{https://doi.org/10.24033/asens.2227}.
\item Supplied literature reference: ``OWR report.'' \url{https://doi.org/10.4171/owr/2019/57}.
\item Literature assessment dated exactly \texttt{2026-08-22}: the supplied curation labels the problem open and states that \(p\)-adic Gross--Zagier and diagonal-cycle results prove important special cases, but not the displayed formula in the extracted generality. This is supplied dated curation derived from the OWR material, not an independently verified literature assessment; the background explicitly requests independent verification of time-sensitive status and cited claims.
\item The supplied statement status is \texttt{corrected\_verified}. Its accompanying qualification says that only formula delimiters, notation, grammar, and typography were repaired, with mathematical content and scope preserved. This label does not establish mathematical correctness or completeness.
\end{itemize}
\end{document}
